% !PRINTED 741
\weberpage{741}
Demnach ist $aM$ und folglich auch $M$ durch $2^{\nu'}$ theilbar,
und wenn wir
\weq{23}{M = 2^{\nu'+\sigma}, \quad N = 2^{\nu'-\sigma-1}}
setzen, so ist $\sigma$ eine nicht negative ganze rationale Zahl.

Setzen wir dies in (16) ein, so folgt:
\[ E = 2^\sigma E' T_0. \]
Nun war in \S.\,197, (6)
\[ E = E' D \]
gesetzt, und es ergiebt sich also:
\[ D = 2^\sigma T_0. \]
Hieraus folgt für den Classenzahlfactor $B$ nach \S.\,199, (15)
und \S.\,200, (21):
\weq{24}{B = 2^{-\sigma} \frac{T}{T_0} = 2^{-\sigma} C,}
worin $C$ eine ungerade ganze Zahl ist.

Demnach ergiebt sich nach \S.\,197, (4) für den zweiten Factor
der Classenzahl $h_1$, der, wie wir gesehen haben, eine \gesperrt{ganze
Zahl} ist:
\weq{25}{h_1 = 2^{-\sigma} h_1' C.}
Nehmen wir nun als bewiesen an, dass $h_1'$ ungerade ist, so
folgt, da auch $C$ ungerade ist, dass $\sigma = 0$ und $h_1$
ungerade ist, und beides ist also damit durch vollständige
Induction bewiesen.

Damit werden die abgeleiteten Resultate folgendermaassen
vereinfacht:

3. \gesperrt{Es ist}
\weq{26}{M = 2^{\nu'}, \quad D = T_0, \quad B = \frac{T}{T_0};}
\gesperrt{der Classenzahlfactor $B$ ist gleich dem Verhältniss
des Regulators $T$ der Einheiten $\tau$ zum Regulator $T_0$
eines Fundamentalsystems von Normaleinheiten, und ist eine
ungerade ganze Zahl.}

Bezeichnen wir die zu $m = 2^\varkappa$ gehörige Zahl $B$ mit
$B_\varkappa$, so ist die Classenzahl $h_1$, wie durch
vollständige Induction folgt,
\weq{27}{h_1 = B_4 B_5 \ldots B_\varkappa,}
und ist also auch eine ungerade Zahl. Damit ist dann, mit
Rücksicht auf \S.\,198, 1., das Haupttheorem bewiesen:

A. \gesperrt{Die Classenzahl im Körper $\Omega_m$ ist, wenn
$m$ eine Potenz von 2 ist, ungerade.}
